\begin{helper}
Fix \(I\subseteq \mathrm{Fin}(n)\), and write \(k=|I|\).  Assume
\(\noderef{ParameterHierarchy}(\varepsilon,\varepsilon_1,\varepsilon_2,n_0)\),
\(n_0<n\),
\[
m=\noderef{TwoBiteNaturalReducedVertexCount}(n),\qquad
p=\noderef{TwoBiteEdgeProbability}\!\left(\frac12,n\right),
\]
and
\[
|I|=k=\noderef{TwoBiteNaturalIndependenceScale}(1+\varepsilon,n).
\]
Let \(\eta,\delta\in\mathbb R\).  Suppose that, for every pair
\(\ell_R,\ell_B\) with \(\ell_R\ell_B\ge k\), the projection-size cell
\[
\noderef{TwoBiteProjectionSizeEvent}(I,k,\ell_R,\ell_B,\omega)
\]
has probability at most
\[
\exp\!\left(
\left(\eta-\frac{2-(\ell_R:\mathbb R)/(k:\mathbb R)
-(\ell_B:\mathbb R)/(k:\mathbb R)}{2}\right)
(k:\mathbb R)\log(n:\mathbb R)\right),
\]
where the displayed quotients and products are interpreted after casting the
natural numbers to real numbers.  Also suppose that the corresponding double
sum
\[
\sum_{\ell_R=0}^{k}\sum_{\ell_B=0}^{k}
\begin{cases}
0,&\ell_R\ell_B<k,\\
\exp\!\left(
\left(\eta-\frac{2-(\ell_R:\mathbb R)/(k:\mathbb R)
-(\ell_B:\mathbb R)/(k:\mathbb R)}{2}\right)(k:\mathbb R)\log(n:\mathbb R)
\right)
\min\!\left(1,
\exp\!\left(-p\bigl(
\noderef{ProjectionOpenPairFunction}(\ell_R,\ell_B,k,p,n)
-\varepsilon^3 k^2\bigr)\right)\right),
&\ell_R\ell_B\ge k
\end{cases}
\]
is at most \(\delta(\binom nk)^{-1}\).  Then
\[
\noderef{TwoBiteEventProbability}\!\left(n,m,p,\,
\{\omega:
I\text{ is independent in }\noderef{TwoBiteFinalGraph}(\omega)
\text{ and }
\noderef{TwoBiteRegularityEvent}(k,\varepsilon,\varepsilon_1,\varepsilon_2,\omega)
\}\right)
\le \delta(\binom nk)^{-1}.
\]
\end{helper}
\begin{proof}
Let
\[
\mathcal B(\omega)\Longleftrightarrow
I\text{ is independent in }\noderef{TwoBiteFinalGraph}(\omega)
\]
and let
\[
\mathcal R(\omega)\Longleftrightarrow
\noderef{TwoBiteRegularityEvent}(k,\varepsilon,\varepsilon_1,\varepsilon_2,\omega).
\]
For \(0\le \ell_R,\ell_B\le k\), put
\[
\mathcal S_{\ell_R,\ell_B}(\omega)\Longleftrightarrow
\noderef{TwoBiteProjectionSizeEvent}(I,k,\ell_R,\ell_B,\omega).
\]
Unfolding \(\noderef{TwoBiteProjectionSizeEvent}\), the cell
\(\mathcal S_{\ell_R,\ell_B}\) asserts three equalities:
\[
|I|=k,\qquad |I.\mathrm{image}(\pi_R^\omega)|=\ell_R,\qquad
|I.\mathrm{image}(\pi_B^\omega)|=\ell_B.
\]
The first equality holds globally by hypothesis.  For every sample
\(\omega\), the two projection-image cardinalities are uniquely determined,
and each is at most \(|I|=k\).  Hence exactly one pair
\((\ell_R,\ell_B)\in\{0,\ldots,k\}^2\) satisfies
\(\mathcal S_{\ell_R,\ell_B}(\omega)\).  Thus these cells are a finite
disjoint partition of the full sample space.

Since \(\noderef{TwoBiteEventProbability}\) is, by definition, the finite sum
of the sample weights over the satisfying samples, finite additivity over this
disjoint partition gives
\[
\Pr(\mathcal B\cap\mathcal R)
=
\sum_{\ell_R=0}^{k}\sum_{\ell_B=0}^{k}
\Pr(\mathcal B\cap\mathcal R\cap\mathcal S_{\ell_R,\ell_B}). \tag{1}
\]

Fix a cell.  If \(\ell_R\ell_B<k\), then the cell is empty.  Indeed, under
\(\mathcal S_{\ell_R,\ell_B}\) the embedded image of \(I\) lies in the
rectangle
\[
I.\mathrm{image}(\pi_R^\omega)\times I.\mathrm{image}(\pi_B^\omega),
\]
whose cardinality is \(\ell_R\ell_B\).  The embedding is injective on \(I\),
so this rectangle must contain the \(k\) distinct embedded points of \(I\),
which is impossible when \(\ell_R\ell_B<k\).  Therefore the corresponding
summand in (1) is \(0\), matching the zero branch of the displayed sum.

Now suppose \(\ell_R\ell_B\ge k\).  If
\(\Pr(\mathcal S_{\ell_R,\ell_B})=0\), then
\(\Pr(\mathcal B\cap\mathcal R\cap\mathcal S_{\ell_R,\ell_B})=0\), because
the numerator event is a subset of the cell; this is again bounded by the
nonnegative displayed cell contribution.  Otherwise, by unfolding
\(\noderef{TwoBiteConditionalProbability}\),
\[
\Pr(\mathcal B\cap\mathcal R\cap\mathcal S_{\ell_R,\ell_B})
=
\Pr(\mathcal S_{\ell_R,\ell_B})\,
\Pr(\mathcal B\cap\mathcal R\mid \mathcal S_{\ell_R,\ell_B}). \tag{2}
\]
The conditional probability factor is at most \(1\).  Also, unfolding
\(\noderef{ParameterHierarchy}\) gives
\[
0\le\varepsilon_1,\qquad \varepsilon_1\le1,\qquad
\noderef{ParameterHierarchyEventualComparisons}
(\varepsilon,\varepsilon_1,\varepsilon_2,n_0).
\]
Thus \(\noderef{FixedSetConditionalIndependenceBound}\), applied to the fixed
set \(I\), this cell \((\ell_R,\ell_B)\), and the present parameter
hypotheses, gives
\[
\Pr(\mathcal B\cap\mathcal R\mid \mathcal S_{\ell_R,\ell_B})
\le
\exp\!\left(-p\bigl(
\noderef{ProjectionOpenPairFunction}(\ell_R,\ell_B,k,p,n)
-\varepsilon^3k^2\bigr)\right).
\]
Combining this with the trivial \(1\)-bound gives the minimum appearing in the
cell summand.

The assumed projection-size cell probability bound gives
\[
\Pr(\mathcal S_{\ell_R,\ell_B})
\le
\exp\!\left(
\left(\eta-\frac{2-(\ell_R:\mathbb R)/(k:\mathbb R)
-(\ell_B:\mathbb R)/(k:\mathbb R)}{2}\right)(k:\mathbb R)\log(n:\mathbb R)
\right).
\]
Multiplying this inequality by the conditional bound from the previous
paragraph and using (2), each nonempty cell contributes at most the nonzero
branch of the displayed double sum.  Together with the empty-cell case,
substitution into (1) gives
\[
\Pr(\mathcal B\cap\mathcal R)
\le
\sum_{\ell_R=0}^{k}\sum_{\ell_B=0}^{k}(\text{displayed cell summand}).
\]
The final assumed double-sum bound is exactly
\[
\sum_{\ell_R=0}^{k}\sum_{\ell_B=0}^{k}(\text{displayed cell summand})
\le \delta(\binom nk)^{-1}.
\]
Transitivity of the two inequalities proves the claimed probability bound.
\end{proof}
